\hnsection{李群—李代数“词典”}

Lie 自 1874 年起在大量论文中发展出的“有限连续群”理论，在他与 F. Engel§ 合作的宏大著作“变换群理论”（{[}4{]}，1888--1893{]}）中得到系统阐述；该理论是第一卷和第三卷最后五章的研究对象，第二卷则专门讨论接触变换。

正如标题所示，这部著作只关注式 (4) 意义下的变换群，其中“变量”空间 \(x_{i}\) 和“参数”空间 \(a_{i}\) 起初扮演着如此重要的角色。另一方面，当时“抽象”群的概念尚未被清楚地分离出来；Lie 在 1885 年（{[}3/j{]}，§5）注意到，在 (10) 的记号下，给出群中两个变换复合参数的方程 \(w = \mathrm{H}(u, v)\) 定义了一个新群，他把它看作参数空间上的变换群，从而得到他所谓的“参数群”（他甚至得到两个，这正是左平移群和右平移群†）。

式 (4) 中的变量 \(x_{i}\) 和参数 \(a_{i}\) 原则上假定为复数（第三卷第 XIX--XXIV 章除外），函数 \(f_{i}\) 假定为解析的；Lie 和 Engel 当然知道，这些函数一般并非对 \(x_{i}\) 和 \(a_{i}\) 的所有复值都有定义，因而这类函数的复合会产生严重困难（{[}4{]}，第 1 卷，第 15--17、33--40 页及各处）；尽管他们在全文中几乎总是写得仿佛所研究的变换可以无条件复合，这无疑是为了论述方便，但在必要时他们明确重申了“局部”观点（例如参见同上第 168 或 189 页，或同书第 3 卷第 2 页页下注）；换言之，他们研究的数学对象接近于本书所谓的运算律片段。他们有时也考虑全局群，例如经典群的 4 个系列（{[}4{]}，第 3 卷，第 682 页），但似乎没有问过一般而言什么构成“全局群”；他们满足于为经典群的“参数”（这些群的“变量”并不造成困难，因为所涉及的变换是 \(\mathbf{C}^n\) 的线性变换）取得恒等变换邻域中的“局部”参数系，而不去考虑所写公式的有效域。不过，他们提出了一个从局部理论自然产生的问题：研究“混合”群，即具有有限个连通分支的群，例如正交群（{[}4{]}，第 1 卷，第 7 页）。他们把这一研究表述为：研究一个在复合及取逆下稳定的变换集合，它是若干集合 \(\mathrm{H}_j\) 的并；每个集合由如 (4) 中的函数系统（\(f_i^{(D)}\)）描述；每个 \(\mathrm{H}_j\) 的（本质）参数个数起初也假定依赖于 \(j\)，但他们证明事实上这个数对所有 \(\mathrm{H}_j\) 都相同。他们的主要结果于是是：存在一个有限连续群 G，使 \(\mathrm{H}_j = \mathrm{G} \circ h_j\) 对某个 \(h_j \in \mathrm{H}_j\) 及所有 \(j\) 成立；他们还证明 G 在混合群中正规，并指出，后者不变量的确定归结为 G 的不变量和一个不连续群的不变量（{[}4{]}，第 1 卷，第 18 章）。

在 {[}4{]} 中发展起来的一般理论最终形成了一部“词典”（作者们没有非常系统地如此表述），把“有限连续群”的性质翻译为其无穷小变换集合的性质。它建立在“Lie 的三个定理”之上，每个定理都由一个断言及其逆命题组成。

第一定理（{[}4{]}，第 1 卷，第 33、72 页以及第 3 卷，第 563 页）首先断言：若 (4) 中的参数是本质的，则函数 \(f_{i}\) 满足如下形式的偏微分方程组

\[
\frac {\partial f _ {i}}{\partial a _ {j}} = \sum_ {k = 1} ^ {r} \xi_ {k i} (f (x, a)) \psi_ {k j} (a) \quad (1 \leqslant i \leqslant n)\tag{15}
\]

其中矩阵 \((\xi_{kl})\) 秩最大且 \(\det(\psi_{kj}) \neq 0\)；反之，若函数 \(f_{i}\) 具有这一性质，则公式 (4) 定义一个变换的群芽。

第二定理（{[}4{]}，第 1 卷，第 149、158 页以及第 3 卷，第 590 页）给出一方面的 \(\xi_{ki}\) 与另一方面的 \(\psi_{ij}\) 之间的关系：关于 \(\xi_{ki}\) 的条件可写成

\[
\sum_ {k = 1} ^ {n} \left(\xi_ {i k} \frac {\partial \xi_ {f l}}{\partial x _ {k}} - \xi_ {j k} \frac {\partial \xi_ {t l}}{\partial x _ {k}}\right) = \sum_ {k = 1} ^ {r} c _ {i j} ^ {k} \xi_ {k l} (1 \leqslant i, j \leqslant r, 1 \leqslant l \leqslant n)\tag{16}
\]

其中 \(c_{ii}^{k}\) 是常数 \((1\leqslant i,j,k\leqslant r)\)，关于 \(i,j\) 斜对称。Maurer 给出的 \(\psi_{ij}\) 条件（{[}10{]}）为：

\[
\frac {\partial \psi_ {k l}}{\partial a _ {m}} - \frac {\partial \psi_ {k m}}{\partial a _ {l}} = \frac {1}{2} \sum_ {1 \leqslant i, j \leqslant r} c _ {i j} ^ {k} (\psi_ {i l} \psi_ {j m} - \psi_ {j l} \psi_ {i m}) (1 \leqslant k, l, m \leqslant r).\tag{17}
\]

引入逆矩阵 \((\alpha_{ij})\)（它是 \((\psi_{ij})\) 的逆矩阵）以及无穷小变换

\[
\mathbf {X} _ {k} = \sum_ {i = 1} ^ {n} \xi_ {k i} \frac {\partial}{\partial x _ {i}}, \quad \mathbf {A} _ {k} = \sum_ {j = 1} ^ {r} \alpha_ {k j} \frac {\partial}{\partial a _ {j}} (1 \leqslant k \leqslant r)\tag{18}
\]

\begin{enumerate}
\def\labelenumi{(\arabic{enumi})}
\setcounter{enumi}{15}
\item
  并且 (17) 分别可写成：
\item
\end{enumerate}

\[
[ \mathrm{X} _ {i}, \mathrm{X} _ {j} ] = \sum_ {k = 1} ^ {r} c _ {i j} ^ {k} \mathrm{X} _ {k} (1 \leqslant i, j \leqslant r).\tag{20}
\]

\[
[ \mathbf {A} _ {i}, \mathbf {A} _ {j} ] = \sum_ {k = 1} ^ {r} c _ {i j} ^ {k} \mathbf {A} _ {k}
\]

反之，若给定 r 个无穷小变换 \(X_{k}\)（\(1 \leqslant k \leqslant r\)），它们线性无关且满足条件 (19)，则由这些变换生成的单参数子群生成一个具有 r 个本质参数的变换群。

最后，第三定理（{[}4{]}，第 1 卷，第 170、297 页以及第 3 卷，第 597 页）把满足 (19) 的无穷小变换系统 \((\mathbf{X}_k)_{1 \leqslant k \leqslant r}\) 的确定归结为一个纯代数问题：必须成立：

\begin{enumerate}
\def\labelenumi{(\arabic{enumi})}
\setcounter{enumi}{20}
\tightlist
\item
\end{enumerate}

\[
c _ {i j} ^ {k} + c _ {j i} ^ {k} = 0\tag{22}
\]

\[
\sum_ {l = 1} ^ {r} \left(c _ {i l} ^ {m} c _ {j k} ^ {l} + c _ {k i} ^ {m} c _ {i j} ^ {l} + c _ {j l} ^ {m} c _ {k l} ^ {l}\right) = 0 \quad (1 \leqslant i, j, k, m \leqslant r).
\]

反之，†若 (21) 和 (22) 成立，则存在一个满足关系 (19) 的无穷小变换系统，从而得到一个具有 r 个参数的变换群（换言之，\(X_{k}\) 的常系数组合构成一个李代数，反过来每个有限维李代数都可由此方式得到）。

这些结果还通过研究同构问题而得到补充。若能通过对变量作可逆坐标变换、对参数作可逆坐标变换而从一个变换群得到另一个，则称两个变换群相似；Lie 从研究一开始就在定义“典范参数”时自然地关注这一概念。他证明，若通过对“变量”作变换可以把一个群的无穷小变换化为另一个群的无穷小变换，则两个群相似（{[}4{]}，第 1 卷，第 329 页）。必要条件是这两个群的李代数同构，Lie 称这些群为“gleichzusammengesetzt”；但这一条件并不充分，因而有一整章（{[}4{]}，第 1 卷，第 19 章）专门讨论保证群“相似”的补充条件。另一方面，置换群理论提供了两个此类群“全同构”的概念（底层“抽象”群的同构）；Lie 将这一概念移植到变换群，并证明两个此类群“全同构”当且仅当它们的李代数同构（{[}4{]}，第 1 卷，第 418 页）。特别地，每个变换群都与它的每个参数群全同构；这说明，当我们希望研究群的结构时，它所作用于其上的“变量”并不重要，实际上重要的只有李代数。†

始终类比于置换群理论，Lie 引入了子群、正规子群和“满同态”（满射同态）的概念，并证明它们分别对应于子代数、理想和满射李代数同态；他此前已经遇到过一个特别重要的“满同态”例子，即伴随表示，并认识到它与群的中心的关系（{[}3f{]}，§ 3，第 9 号）。对于这些结果以及基本定理，关键工具是 Jacobi--Clebsch 定理，它给出了微分系统的完全可积性（这种定理的形式之一称为“{[}Frobenius 定理”）；不过，他使用单参数群给出了新的证明（{[}4{]}，第 1 卷，第 6 章）。

对置换群如此重要的传递性和本原性概念，同样自然地出现在“有限连续”变换群中，Lie--Engel 著作对此作了详细研究（{[}4{]}，第 1 卷，第 13 章及各处）；人们也认识到了它们与一点的稳定子群以及齐性空间概念之间的关系（在不采取全局观点所可能的范围内）（{[}4{]}，第 1 卷，第 425 页）。

最后，在 {[}4{]} 中，这部“词典”通过引入导出群和可解群的概念而得以完成（Lie 称之为“可积分群”；这一术语由微分方程理论启发，一直沿用到 H. Weyl 的著作）（{[}4{]}，第 1 卷，第 261 页及第 3 卷，第 678--679 页）；Lie 还在别处于 1883 年认识到了换位子与括号之间的关系（{[}3{]}，第 5 卷，第 358 页）。

